\chapter{AgenTwin Implementation}
\label{chap:implementation}

This chapter describes the AgenTwin proof-of-concept implementation,
    realizing the architecture presented in Chapter~\ref{chap:architecture}
    through a concrete deployment integrating heterogeneous Digital Twin platforms
    mediated by AI agents.
Figure~\ref{fig:poc-architecture} illustrates the implementation architecture,
    showing the integration of KTWIN and Eclipse Ditto platforms through agent-based coordination.

\begin{figure}[htbp]
    \centering
    \includegraphics[width=\textwidth]{figures/proof-of-concept.pdf}
    \caption{AgenTwin proof-of-concept implementation architecture}
    \label{fig:poc-architecture}
\end{figure}

\section{User Application Component}
\label{sec:user-application-component}

The User Application Component provides the end-user interface through which
    requests are submitted to the Digital Ecosystem.
To ensure research reproducibility while avoiding proprietary constraints,
    we use \textbf{Open WebUI}~\cite{baek2025},
    an open-source self-hosted AI interface that enables flexible agent protocol integration.

Open WebUI serves as both the user interface and agent client,
    translating natural language user requests into structured agent communications.
The platform's extensibility allows implementation of custom Agent-to-Agent (A2A) protocol clients
    without vendor lock-in to proprietary chat interfaces or cloud services.

\section{Digital Ecosystem Orchestration}
\label{sec:digital-ecosystem-orchestration}

The Digital Ecosystem's agent communication is realized through
    \textbf{Microsoft Agent Framework}.
The framework's Workflow layer provides the graph-based multi-agent orchestration
    underpinning the Group Chat Manager component described in
    Section~\ref{subsec:session-layer},
    allowing participating agents to share conversation context and
    interact without predetermined communication sequences.

Agent Framework's model-agnostic design enables evaluation across both
    cloud-hosted backends (e.g.\ Claude Sonnet 3.5 via Anthropic API) and
    locally-hosted alternatives (e.g.\ Llama\,3 via Ollama).
This facilitates prototyping under realistic constraints while
    supporting performance comparison and cost-sensitive deployment scenarios.

\section{Digital Twin Platform Components}
\label{sec:digital-twin-platform-components}

Two open-source DT platforms were selected to demonstrate heterogeneous system integration
    while maintaining realistic semantic complexity.
The selection criteria prioritized platforms with
    sufficient heterogeneity to require semantic mediation, and
    extensibility to support MCP server integration.

\textbf{KTWIN}~\cite{Wermann_Wickboldt_2025} implements the Smart City DT platform.
KTWIN is a Kubernetes-based serverless framework designed to 
    simplify DT adoption through cloud-native deployment patterns.
It natively supports the Digital Twins Definition Language (DTDL) and has been
    validated for smart city applications using the
    \texttt{opendigitaltwins-smartcities} ontology~\cite{russomAzureOpendigitaltwinssmartcities2021}.
This ontology provides comprehensive models for
    urban infrastructure entities including
    parking facilities,
    EV charging stations,
    traffic sensors, and
    building management systems.

Key technical characteristics include:

\begin{itemize}
    \item \textbf{Serverless Event-Driven Architecture:}
    Function-as-a-Service (FaaS) deployment enables
        automatic scaling and
        reduced operational overhead, while
        the centralized event broker (RabbitMQ) supports
        real-time state updates through
        multiple protocols (MQTT, HTTP, AMQP).

    \item \textbf{Kubernetes-Native Infrastructure:}
    Custom Resource Definitions (CRDs) and operators automate
        infrastructure provisioning based on
        DTDL ontology definitions,
        enabling multi-cloud and edge deployments without vendor lock-in.
    
    \item \textbf{Ontology-Driven Automation:}
    Native DTDL support eliminates schema translation overhead,
        allowing direct implementation of the \texttt{smartcities} ontology concepts.
\end{itemize}

\textbf{Eclipse Ditto} implements the Energy Grid DT platform.
Ditto is a mature open-source IoT platform providing
    DT capabilities through its ``Thing'' abstraction,
    widely adopted in industrial IoT deployments.
For this research, Ditto manages twins conforming to the
    \texttt{opendigitaltwins-energygrid} ontology~\cite{raviAzureOpendigitaltwinsenergygrid2021}, which adapts the
    Common Information Model (CIM) standard---the energy industry's reference ontology---to DTDL.
Since Ditto's native knowledge representation uses the Thing model,
    schema translation will be necessary for this research.

Key technical characteristics include:

\begin{itemize}
    \item \textbf{Thing-Based Knowledge Representation:}
    Distinct from KTWIN's DTDL approach,
        necessitating genuine semantic mapping rather than trivial schema alignment
    
    \item \textbf{Rich REST and WebSocket APIs:}
    Well-documented interfaces simplify MCP server implementation for resource access
    
    \item \textbf{RabbitMQ Integration:}
    Similar messaging infrastructure to KTWIN reduces deployment complexity
        while maintaining semantic heterogeneity at the ontology level
    
    \item \textbf{Production-Grade Maturity:}
    Extensive documentation, active community, and industrial adoption
        ensure realistic platform behavior for evaluation
\end{itemize}

The deliberate selection of platforms with similar
    messaging infrastructure (RabbitMQ) but
    different knowledge representations (native DTDL vs. Thing model)
    creates the realistic heterogeneity necessary to
    validate agent-mediated semantic interoperability while
    maintaining implementation feasibility within research constraints.
Neither platform provides native MCP server capabilities.
Implementing MCP servers for both KTWIN and Eclipse Ditto
    constitutes part of this research scope.

\section{Ontology Preparation}
\label{sec:ontology-preparation}

Designing the cross-domain use case revealed an immediate representational problem:
    the geographic correlation between Smart City charging station locations and
    Energy Grid substation service zones cannot
    be represented in the Energy Grid graph because the
    ontology does not expose a \texttt{Location} relationship at the equipment level.
Without this association, it's impossible to instantiate the structural semantic gap
    described in Section~\ref{subsec:semantic-gaps}.

The published ontology uses \texttt{SubGeographicalRegion} to
    represent areas such as cities and districts.
Within these regions, \texttt{Substation} describe a collection of equipments
    like \texttt{PowerTransformer} through which electric energy in bulk is passed for the purposes of
    switching or modifying its characteristics.
An \texttt{EnergyConsumer} is a generic user of energy, which can be a 
    \texttt{CustomerLoad} (a specific load entity) or an
    \texttt{EquivalentLoad} (a virtual load representing aggregated consumption).
Both transformers and loads are connected through \texttt{Terminal} entities,
    which are in turn connected to \texttt{ConnectivityNode} entities that
    represent physical connection points.

The Smart City graph models the urban environment through
    \texttt{District}, and \texttt{Neighbourhood} entities,
    which are linked by GeoSPARQL spatial containment relationships.
The \texttt{EVChargingStation} entity represents an EV charging infrastructure asset
    located within the Smart City graph.
Besides its administrative area, a charging station spatial representation
    is described by a \texttt{GeoLocation} entity.

Figure~\ref{fig:preliminary-semantic-gap} illustrates this semantic challenge.
Although a substation can be correlated with a neighbourhood through spatial containment,
    it is not possible to correlate a charging station to its customer load counterpart because
    the load-level entities have no geographic grounding.

\begin{figure}[htbp]
    \centering
    \includegraphics[height=0.85\textheight,keepaspectratio]{figures/preliminary-semantic-gap.pdf}
    \caption{Cross-domain ontology model before the patch}
    \label{fig:preliminary-semantic-gap}
\end{figure}

A systematic review of all CIM associations that reference \texttt{Location}
    was conducted to determine whether any standard-conformant path could
    bring geographic grounding to the load-level entities.
The answer was found in \texttt{PowerSystemResource}.
Its location is normatively defined in
    IEC~61968-13:2008~\cite{internationalelectrotechnicalcommissionApplicationIntegrationElectric2008},
    the CIM RDF exchange format for distribution networks,
    whose CDPSM profile uses the association throughout its RDF examples
    for every equipment class and explicitly lists \texttt{Location} as a class used in CDPSM
    but not in CPSM (the transmission-only profile~\cite{Uslar_2012}).
The Energy Grid repository describes its models as adapted from
    the CIM standards suite (via the CIM Users Group) and from the
    \texttt{smart-data-models/dataModel.EnergyCIM} data models~\cite{abellaSmartdatamodelsDataModelEnergyCIM2024},
    the latter of which already includes an explicit \texttt{Location} property in its
    \texttt{PowerSystemResource} schema.
Its omission from the DTDL ontology therefore reflects an incomplete adaptation,
    which the repository README acknowledges by describing the published models as
    a \textit{``first iteration''} and explicitly inviting contributors to
    add missing entities before the next release.

A patch was applied to a fork of the ontology~\cite{correaHenriKCorreaOpendigitaltwinsenergygrid2026}:
    the \texttt{Location} relationship was restored to \texttt{PowerSystemResource},
    making the Energy Grid graph traversable by geographic position.
The restored \texttt{Location} is a structured CIM entity comprising
    a sequence of \texttt{PositionPoints} together with a \texttt{CoordinateSystem} relationship
    whose \texttt{crsUrn} property holds an OGC URN declaring the coordinate reference system.
All other design decisions were retained verbatim as
    authentic representations of the published ontology.

One admissibility criterion governed this decision:
    a modification is admissible only if it
    (a)~restores a property normatively defined in the upstream standard,
    (b)~is confirmed by at least one independent CIM implementation, and
    (c)~increases rather than reduces the semantic complexity the agent must resolve at runtime.
The \texttt{Location} patch satisfies all three conditions.
The motivation was to make the research problem well-formed;
    the effect was to introduce additional semantic complexity.

\section{Agent-Platform Integration}
\label{sec:agent-interaction-impl}

While Section~\ref{subsec:session-layer} described
    agent-to-agent coordination through the A2A protocol,
    this section focuses on the complementary agent-to-platform interaction pattern
    realized through the Model Context Protocol (MCP).
MCP enables dynamic resource discovery and query execution without requiring agents to
    maintain platform-specific client code or
    hardcoded schema knowledge.

Each DT platform exposes an MCP server.
The server publishes three types of information:
    resource catalogs (available entity types and instances),
    schema metadata (properties, relationships, semantic annotations), and
    query capabilities (supported operations like filtering, spatial queries, aggregation).
For KTWIN, the MCP server wraps Kubernetes CRDs and RabbitMQ event streams,
    exposing \texttt{smart-city} ontology entities.
For Ditto, the MCP server interfaces with REST and WebSocket APIs,
    querying \texttt{energy-grid} ontology entities.

Each domain agent embeds an MCP client
    that connects to its associated platform's MCP server.
The client provides the agent's LLM with tool definitions for
    discovering resources,
    retrieving schemas, and
    executing queries with filtering and projection.
This architecture enables agents to dynamically discover platform capabilities through MCP,
    negotiate semantic mappings through A2A conversations, and
    execute coordinated queries to solve cross-domain problems---all
    without requiring modifications to the underlying DT platforms beyond MCP server deployment.

\section{Task Completion Contract}
\label{sec:task-completion-contract}
% TODO: create figure (mermaid) illustrating the task completion contract mechanism
Each A2A task submitted to the Digital Ecosystem includes a terminal tool schema
    encoded as an AgenTwin-defined key in the A2A message metadata,
    which the protocol defines as a flexible key-value map for deployment-specific context.
The schema is defined by the Client Agent (the caller) and describes the structured
    output the ecosystem must produce to conclude the task.
The Group Chat Manager receives the schema at task initialisation,
    registers the terminal tool with Agent Framework's plugin registry, % TODO: Explain MAF plugin registry in the background chapter
    and ends the workflow when the tool is invoked by a domain agent.
The tool arguments are captured verbatim and returned to the Client Agent
    as the A2A task artifact.

This design is use-case-agnostic: the Group Chat Manager contains no knowledge of
    domain entities, or any specific output schema.
Any deployment that needs a different output structure provides its own
    terminal tool name and JSON Schema in the A2A task metadata,
    and the ecosystem adapts without modification.

The current proof-of-concept operates under a closed trust model:
    the Group Chat Manager accepts the terminal tool schema from the Client Agent
    without authenticating the caller or validating the schema against
    an authorisation policy.
Production deployments at A2A trust boundaries would require
    mutual authentication and schema governance
    before allowing an external caller to register executable tools.

This mechanism bridges Agent Framework's plugin registration model
    (Section~\ref{subsec:ai-agent-frameworks}) with A2A's agent-driven task completion model
    (Section~\ref{subsec:ai-agent-protocols}), without modifying either framework.
